\begingroup
\begin{table*}[!t]
    \centering
    \small
    \scalebox{0.9}{
    \begin{tabular}{p{\linewidth}}
        \toprule
        \underline{\textbf{\textsc{Exemplar for Date Understanding}}} \\
        \vspace{-2mm}
        \# Q: Yesterday was April 30, 2021. What is the date tomorrow in MM/DD/YYYY? \\
        \vspace{-1mm}
        \# To answer this question, we write a program to answer the following subquestions: \\
        \# import relevant packages \\
        from datetime import date, time, datetime \\
        from dateutil.relativedelta import relativedelta \\
        \# 1. What is the date yesterday? (independent, support: ["Yesterday was April 30, 2021"]) \\
        date\_yesterday = date(2021,4,30) \\
        \# 2. What is the date today? (depends on 1, support: ["Yesterday was April 30, 2021"]) \\
        date\_today = date\_yesterday + relativedelta(days=1) \\
        \# 3. What is the date tomorrow? (depends on 2, support: []) \\
        date\_tomorrow = date\_today + relativedelta(days=1) \\
        \# 4. Final Answer: What is the date tomorrow in MM/DD/YYYY? (depends on 3, support: []) \\
        answer = date\_tomorrow.strftime("\%m/\%d/\%Y") \\
        \bottomrule
    \end{tabular}
    }
        \caption{
    An exemplar from our prompt for Date Understanding.
    }
        \label{tab:appendix-date_prompt}

\end{table*}

\endgroup